% Generated draft: requires mathematical and source review.
\documentclass[11pt]{article}
\usepackage[margin=25mm]{geometry}
\usepackage{fontspec}
\setmainfont{Latin Modern Roman}
\usepackage{amsmath,amssymb,mathtools,unicode-math}
\setmathfont{Latin Modern Math}
\usepackage{enumitem,needspace,xurl,hyperref}
\setlength{\parindent}{0pt}
\setlength{\parskip}{5pt}
\setlength{\emergencystretch}{3em}
\newcommand{\N}{\mathbb{N}}
\newcommand{\Z}{\mathbb{Z}}
\newcommand{\Q}{\mathbb{Q}}
\newcommand{\R}{\mathbb{R}}
\newcommand{\C}{\mathbb{C}}
\newcommand{\F}{\mathbb{F}}
\newcommand{\A}{\mathbb{A}}
\newcommand{\PP}{\mathbb{P}}
\newcommand{\E}{\mathbb{E}}
\newcommand{\eps}{\varepsilon}
\newcommand{\dd}{\,\mathrm{d}}
\begin{document}
\section*{UnsolvedMath ID: 30000290}
\textbf{Source problem code:} catalogue entry
\subsection{Definitions and mathematical statement}
Let $T$ be a continuous linear operator on the topological vector space $X$ intended by the report. For a set $A\subseteq\N_0$, its lower asymptotic density is
\[
 \underline{\operatorname{dens}}(A)
 =\liminf_{m\to\infty}\frac{\#(A\cap\{0,\ldots,m-1\})}{m}.
\]
A vector $x\in X$ is frequently hypercyclic for $T$ if, for every nonempty open set $U\subseteq X$,
\[
 \underline{\operatorname{dens}}\{n\in\N_0:T^nx\in U\}>0.
\]
Write $\operatorname{FHC}(T)$ for this set; $T$ is frequently hypercyclic when $\operatorname{FHC}(T)\ne\varnothing$.

\begin{enumerate}
\item Characterize frequently hypercyclic weighted backward shifts. For a unilateral sequence space with coordinates indexed by $\N_0$, a conventional indexing is
\[
 (B_wx)_n=w_{n+1}x_{n+1}\qquad(n\ge0),
\]
with weights chosen so that $B_w$ is a well-defined continuous operator on that space. A bilateral version uses $n\in\Z$. The source does not say which index set, sequence spaces, scalar field, or weight hypotheses belong to the requested characterization, so neither convention is selected as the full target.
\item Can $\operatorname{FHC}(T)$ be of second category in $X$ for an operator in the intended class? Here second category means nonmeagre: the set is not a countable union of nowhere dense subsets of $X$. This is an existence question and is not replaced by the stronger claim that $\operatorname{FHC}(T)$ is comeagre. The source does not restrict this second question to weighted shifts.
\end{enumerate}

The supplied dated assessment says that the first question is solved for weighted shifts on $\ell^p$ by the Bayart--Ruzsa criterion, while it labels the bundled second-category question open in its stated generality. It provides neither the precise range of $p$ and shift conventions nor the criterion itself. Those claims are retained as dated curation and must be checked, rather than promoted to a current theorem or status assessment. The operator and ambient-space class for the second question also remain unspecified.
\subsection{Short English statement}
One question asks which weighted backward shifts visit every open set with positive lower frequency. The other asks whether the collection of such vectors for an operator can be nonmeagre. Their scopes and the exact hypotheses of the cited characterization require source review.
\subsection{Sources}
\begin{itemize}
\item Local UnsolvedMath registry, permanent ID 30000290, code catalogue entry, ``Frequently Hypercyclic Weighted Backward Shifts.'' Proposed year 2005; proposer unspecified.
\item \emph{Mini-Workshop: Operators on Spaces of Analytic Functions} (2006), Oberwolfach Reports, pp.\ 2043--2090, DOI 10.4171/owr/2005/36; \url{https://doi.org/10.4171/owr/2005/36}. The displayed year 2006 differs from the DOI and proposed year 2005.
\item Supplied reference labeled ``Bayart--Ruzsa characterization'': \url{https://doi.org/10.1017/S0143385713000771}.
\item Statement status corrected\_verified, with a supplied claim that the official report was checked and wording made self-contained with normalized notation.
\item Literature triage checked 2026-08-21: partially solved. It describes the $\ell^p$ characterization as solved and the second-category question as open in this generality. The full hypotheses and up-to-date accuracy require independent verification. Created 2026-08-25; updated 2026-09-16.
\item Registry curation by Ulam AI / UnsolvedMath Contributors, CC BY 4.0; underlying report and literature provenance remain as supplied. Unreviewed statement-only draft.
\end{itemize}
\end{document}
